% GENERATED FILE. Edit canonical lessons, then run npm run generate:books.
% canonical-source-hash: fnv1a64:468ae89c2732e06d
% canonical-lessons: BN-C166-shokun, BN-C166-timi, BN-C166-kakra, BN-C166-shamuk, BN-C166-kriket

\chapter{Vulture, Whale, Crab, Snail, Cricket}
\label{ch:bn-vulture-whale-crab-snail-cricket}
\hlchaptermodality{166}

\begin{chapteropening}
\textbf{By the end of this chapter:} \emph{I can say a vulture, a whale, a crab, a snail and cricket.}

Complete the last lesson of chapter 166: I can say a vulture, a whale, a crab, a snail and cricket.
\end{chapteropening}

\section[\texorpdfstring{\bn{শকুন} (shôkun)}{শকুন (shôkun)}]{\bn{শকুন} (shôkun) — a vulture}

\label{lesson:BN-C166-shokun}



\begin{quote}
Before the new one: say the Bengali for sandesh, then the Bengali for a paratha.
\end{quote}

\subsection*{You'll want to know: \bn{শকুন}}
\textbf{\bn{শকুন}} — \emph{shôkun} — \textquotedblleft{}a vulture\textquotedblright{}.

\begin{cousinweb}
\textbf{\bn{শকুন}} (\emph{shôkun}) is Sanskrit \emph{śakuna}, \textquotedblleft{}bird\textquotedblright{}, taken over whole as a learned loan, a tatsama (CDIAL 12245). In Bengali it has narrowed to \textquotedblleft{}vulture\textquotedblright{}.
\end{cousinweb}

\begin{grammarlens}[title={What you've built}]
One more word for this chapter.
\end{grammarlens}

\subsection*{Your turn}
\begin{itemize}
\raggedright
  \item \emph{Say it:} \emph{shôkun}
  \item \emph{Say it:} \emph{shôkun}, once more
  \item \emph{Say it:} say \emph{shôkun}
  \item \emph{Recall:} say the Bengali for sandesh, then the Bengali for a paratha, then say \emph{shôkun} again
\end{itemize}

\begin{tcolorbox}[breakable,colback=teal!4,colframe=teal!35!black,title={Before you move on}]
What does \bn{শকুন} mean? (\textquotedblleft{}A vulture\textquotedblright{}.) Say it once more.
\end{tcolorbox}
\section[\texorpdfstring{\bn{তিমি} (timi)}{তিমি (timi)}]{\bn{তিমি} (timi) — a whale}

\label{lesson:BN-C166-timi}



\begin{quote}
Before the new one: say the Bengali for a paratha, then the Bengali for a vulture.
\end{quote}

\subsection*{You'll want to know: \bn{তিমি}}
\textbf{\bn{তিমি}} — \emph{timi} — \textquotedblleft{}a whale\textquotedblright{}.

\begin{cousinweb}
\textbf{\bn{তিমি}} (\emph{timi}) is Sanskrit \emph{timi}, \textquotedblleft{}a large fabled fish\textquotedblright{}, taken over whole as a learned loan, a tatsama.
\end{cousinweb}

\begin{grammarlens}[title={What you've built}]
One more word for this chapter.
\end{grammarlens}

\subsection*{Your turn}
\begin{itemize}
\raggedright
  \item \emph{Say it:} \emph{timi}
  \item \emph{Say it:} \emph{timi}, once more
  \item \emph{Say it:} say \emph{timi}
  \item \emph{Recall:} say the Bengali for a paratha, then the Bengali for a vulture, then say \emph{timi} again
\end{itemize}

\begin{tcolorbox}[breakable,colback=teal!4,colframe=teal!35!black,title={Before you move on}]
What does \bn{তিমি} mean? (\textquotedblleft{}A whale\textquotedblright{}.) Say it once more.
\end{tcolorbox}
\section[\texorpdfstring{\bn{কাঁকড়া} (kā̃kṛā)}{কাঁকড়া (kā̃kṛā)}]{\bn{কাঁকড়া} (kā̃kṛā) — a crab}

\label{lesson:BN-C166-kakra}



\begin{quote}
Before the new one: say the Bengali for a vulture, then the Bengali for a whale.
\end{quote}

\subsection*{You'll want to know: \bn{কাঁকড়া}}
\textbf{\bn{কাঁকড়া}} — \emph{kā̃kṛā} — \textquotedblleft{}a crab\textquotedblright{}.

\begin{cousinweb}
\textbf{\bn{কাঁকড়া}} (\emph{kāṁkṛā}) is a tadbhava, a word inherited from Sanskrit and reshaped by sound change. It comes from \emph{karkaṭa}, \textquotedblleft{}crab\textquotedblright{} (CDIAL 2816). CDIAL lists Bengali \emph{kāṁkṛā}, \textquotedblleft{}crab\textquotedblright{}, and suspects that the Sanskrit word is itself of non-Aryan origin.
\end{cousinweb}

\begin{grammarlens}[title={What you've built}]
One more word for this chapter.
\end{grammarlens}

\subsection*{Your turn}
\begin{itemize}
\raggedright
  \item \emph{Say it:} \emph{kā̃kṛā}
  \item \emph{Say it:} \emph{kā̃kṛā}, once more
  \item \emph{Say it:} say \emph{kā̃kṛā}
  \item \emph{Recall:} say the Bengali for a vulture, then the Bengali for a whale, then say \emph{kā̃kṛā} again
\end{itemize}

\begin{tcolorbox}[breakable,colback=teal!4,colframe=teal!35!black,title={Before you move on}]
What does \bn{কাঁকড়া} mean? (\textquotedblleft{}A crab\textquotedblright{}.) Say it once more.
\end{tcolorbox}
\section[\texorpdfstring{\bn{শামুক} (shāmuk)}{শামুক (shāmuk)}]{\bn{শামুক} (shāmuk) — a snail}

\label{lesson:BN-C166-shamuk}



\begin{quote}
Before the new one: say the Bengali for a whale, then the Bengali for a crab.
\end{quote}

\subsection*{You'll want to know: \bn{শামুক}}
\textbf{\bn{শামুক}} — \emph{shāmuk} — \textquotedblleft{}a snail\textquotedblright{}.

\begin{cousinweb}
\textbf{\bn{শামুক}} (\emph{shāmuk}) is a tadbhava, a word inherited from Sanskrit and reshaped by sound change. It comes from \emph{śambūka}, \textquotedblleft{}bivalve shell\textquotedblright{} (CDIAL 12316). CDIAL lists Bengali \emph{sāmuk}, \textquotedblleft{}bivalve shell, snail\textquotedblright{}.
\end{cousinweb}

\begin{grammarlens}[title={What you've built}]
One more word for this chapter.
\end{grammarlens}

\subsection*{Your turn}
\begin{itemize}
\raggedright
  \item \emph{Say it:} \emph{shāmuk}
  \item \emph{Say it:} \emph{shāmuk}, once more
  \item \emph{Say it:} say \emph{shāmuk}
  \item \emph{Recall:} say the Bengali for a whale, then the Bengali for a crab, then say \emph{shāmuk} again
\end{itemize}

\begin{tcolorbox}[breakable,colback=teal!4,colframe=teal!35!black,title={Before you move on}]
What does \bn{শামুক} mean? (\textquotedblleft{}A snail\textquotedblright{}.) Say it once more.
\end{tcolorbox}
\section[\texorpdfstring{\bn{ক্রিকেট} (krikeṭ)}{ক্রিকেট (krikeṭ)}]{\bn{ক্রিকেট} (krikeṭ) — cricket (the game)}

\label{lesson:BN-C166-kriket}



\begin{quote}
Before the new one: say the Bengali for a crab, then the Bengali for a snail.
\end{quote}

\subsection*{You'll want to know: \bn{ক্রিকেট}}
\textbf{\bn{ক্রিকেট}} — \emph{krikeṭ} — \textquotedblleft{}cricket (the game)\textquotedblright{}.

\begin{cousinweb}
\textbf{\bn{ক্রিকেট}} (\emph{krikeṭ}) is a loan from English \emph{cricket}.
\end{cousinweb}

\begin{grammarlens}[title={What you've built}]
One more word for this chapter.
\end{grammarlens}

\subsection*{Your turn}
\begin{itemize}
\raggedright
  \item \emph{Say it:} \emph{krikeṭ}
  \item \emph{Say it:} \emph{krikeṭ}, once more
  \item \emph{Say it:} say \emph{krikeṭ}
  \item \emph{Recall:} say the Bengali for a crab, then the Bengali for a snail, then say \emph{krikeṭ} again
\end{itemize}

\begin{tcolorbox}[breakable,colback=teal!4,colframe=teal!35!black,title={Before you move on}]
What does \bn{ক্রিকেট} mean? (\textquotedblleft{}Cricket (the game)\textquotedblright{}.) Say it once more.
\end{tcolorbox}
